\chapter{全文总结与展望}

\section{全文总结}

本文围绕无人机搭载电磁超声探头进行自主无损检测的核心挑战——探头与被测表面接触状态的可靠判定——开展了从物理建模、平台搭建、信号采集、特征融合到闭环决策的系统性研究。主要工作总结如下：

（1）在理论层面，从电磁-弹性耦合第一性原理出发，将EMAT信号链表示为激励-传播-反射-接收四环节串联的可观测动态过程，推导了洛伦兹体力与弹性动力学方程的基本关系。通过对接触边界和升距效应的回波统计建模，将接触状态从确定性二值标签提升为多因素条件概率表达。

（2）在平台与数据层面，完成了无人机多模态实验平台的全面搭建，集成了Livox MID-360 LiDAR、Intel RealSense D435 RGBD相机、PX4飞控与自研EMAT探头，开发了完整的ROS驱动节点（C++17）和多模态同步数据采集系统。

（3）在算法层面，针对多源异步采样问题，提出了统一时间基准加权插值对齐和模态独立编码器映射方案；设计了物理约束注意力机制，将时间邻近、空间邻近和法向一致性先验以对数偏置形式嵌入Transformer注意力计算中；引入了模态可靠性门控权重和双阈值持续时间约束滞回状态机，形成了从不确定性评估到控制闭环判据的完整证据链。

（4）设计了系统性的实验验证方案和消融实验，为后续完整评估各模块功能贡献和闭环飞行验证提供了明确的实验路线。

\textit{注：本章结论将根据实验完成后获得的定量结果进行更新和具体化。}

\section{后续工作展望}

在本文研究工作的基础上，仍有以下方向值得进一步深入研究：

（1）\textbf{多目标多缺陷类型扩展}：本文目前仅针对平面铝板表面的接触状态进行判别，后续可扩展至曲面、粗糙表面以及含人工缺陷（如裂纹、腐蚀坑）的复杂结构，探讨缺陷类型对接触判定和超声信号的影响。

（2）\textbf{高动态环境迁移}：本文实验在室内受控条件下完成，后续需在室外真实风场、变光照和非结构化表面等复杂环境下验证模型的鲁棒性和泛化能力，可能涉及域自适应和数据增强策略。

（3）\textbf{轻量化与边缘部署}：当前模型以推理精度和物理可解释性为优先目标，后续可将物理约束注意力蒸馏到轻量化架构（如CNN-GRU或MobileNet变体）中，满足NVIDIA Jetson平台上的实时推理需求。

（4）\textbf{接触力-信号质量联合优化}：将本文的接触概率估计与主动阻抗控制相结合，实现对接触力和超声信号质量的联合在线优化，而非当前的开环PID接触维持策略。

（5）\textbf{时序标注自动化}：当前弱标签+人工复核的标注策略耗时较长，后续可探索基于自监督预训练和少量精确标签的半监督学习策略，降低标注成本。